% Ngân hàng câu hỏi TOÁN THỰC TẾ T12 -- Chương 4, Bài 13, Dạng D75
%% Dạng: Ứng dụng hình học của tích phân: thể tích vật thể
%% 100 câu: NB 30 -- TH 40 -- VD 20 -- VDC 10 ; mỗi mức: trắc nghiệm : đúng/sai : trả lời ngắn = 3 : 1 : 1
%% Phương án ĐẦU TIÊN trong \cauTN luôn là đáp án đúng (gói trondethi tự đảo).
%% Phần II: \y{nội dung}{d|s}. Phần III: đáp án đặt ở đối số cuối.

\cauTN[id=T12-C4-B13-D75-NB-001,muc=NB,loigiai={$V=\pi\int_0^{9}\left(\sqrt{3x}\right)^2dx=\pi\int_0^{9}3x\,dx=\pi\left[\dfrac{3x^2}{2}\right]_0^{9}=\dfrac{243\pi}{2}\approx381{,}7$ m$^3$.}]{Một bồn chứa được tạo thành khi quay quanh trục hoành hình phẳng giới hạn bởi $y=\sqrt{3x}$, $y=0$, $x=0$ và $x=9$ (đơn vị trên trục là mét). Thể tích của bồn chứa bằng bao nhiêu mét khối (làm tròn đến hàng phần trăm)?}
  {$381{,}7$.}
  {$763{,}41$.}{$121{,}5$.}{$190{,}85$.}

\cauTN[id=T12-C4-B13-D75-NB-002,muc=NB,loigiai={$V=\pi\int_0^{8}\left(\sqrt{x}\right)^2dx=\pi\int_0^{8}1x\,dx=\pi\left[\dfrac{x^2}{2}\right]_0^{8}=\dfrac{64\pi}{2}\approx100{,}53$ m$^3$.}]{Một bồn chứa được tạo thành khi quay quanh trục hoành hình phẳng giới hạn bởi $y=\sqrt{x}$, $y=0$, $x=0$ và $x=8$ (đơn vị trên trục là mét). Thể tích của bồn chứa bằng bao nhiêu mét khối (làm tròn đến hàng phần trăm)?}
  {$100{,}53$.}
  {$201{,}06$.}{$32$.}{$50{,}27$.}

\cauTN[id=T12-C4-B13-D75-NB-003,muc=NB,loigiai={$V=\pi\int_0^{3}\left(\sqrt{x}\right)^2dx=\pi\int_0^{3}1x\,dx=\pi\left[\dfrac{x^2}{2}\right]_0^{3}=\dfrac{9\pi}{2}\approx14{,}14$ m$^3$.}]{Một bồn chứa được tạo thành khi quay quanh trục hoành hình phẳng giới hạn bởi $y=\sqrt{x}$, $y=0$, $x=0$ và $x=3$ (đơn vị trên trục là mét). Thể tích của bồn chứa bằng bao nhiêu mét khối (làm tròn đến hàng phần trăm)?}
  {$14{,}14$.}
  {$28{,}27$.}{$4{,}5$.}{$7{,}07$.}

\cauTN[id=T12-C4-B13-D75-NB-004,muc=NB,loigiai={$V=\pi\int_0^{4}\left(\sqrt{3x}\right)^2dx=\pi\int_0^{4}3x\,dx=\pi\left[\dfrac{3x^2}{2}\right]_0^{4}=\dfrac{48\pi}{2}\approx75{,}4$ m$^3$.}]{Một bồn chứa được tạo thành khi quay quanh trục hoành hình phẳng giới hạn bởi $y=\sqrt{3x}$, $y=0$, $x=0$ và $x=4$ (đơn vị trên trục là mét). Thể tích của bồn chứa bằng bao nhiêu mét khối (làm tròn đến hàng phần trăm)?}
  {$75{,}4$.}
  {$150{,}8$.}{$24$.}{$37{,}7$.}

\cauTN[id=T12-C4-B13-D75-NB-005,muc=NB,loigiai={$V=\pi\int_0^{7}\left(\sqrt{4x}\right)^2dx=\pi\int_0^{7}4x\,dx=\pi\left[\dfrac{4x^2}{2}\right]_0^{7}=\dfrac{196\pi}{2}\approx307{,}88$ m$^3$.}]{Một bồn chứa được tạo thành khi quay quanh trục hoành hình phẳng giới hạn bởi $y=\sqrt{4x}$, $y=0$, $x=0$ và $x=7$ (đơn vị trên trục là mét). Thể tích của bồn chứa bằng bao nhiêu mét khối (làm tròn đến hàng phần trăm)?}
  {$307{,}88$.}
  {$615{,}75$.}{$98$.}{$153{,}94$.}

\cauTN[id=T12-C4-B13-D75-NB-006,muc=NB,loigiai={$V=\pi\int_0^{6}\left(\sqrt{3x}\right)^2dx=\pi\int_0^{6}3x\,dx=\pi\left[\dfrac{3x^2}{2}\right]_0^{6}=\dfrac{108\pi}{2}\approx169{,}65$ m$^3$.}]{Một bồn chứa được tạo thành khi quay quanh trục hoành hình phẳng giới hạn bởi $y=\sqrt{3x}$, $y=0$, $x=0$ và $x=6$ (đơn vị trên trục là mét). Thể tích của bồn chứa bằng bao nhiêu mét khối (làm tròn đến hàng phần trăm)?}
  {$169{,}65$.}
  {$339{,}29$.}{$54$.}{$84{,}82$.}

\cauTN[id=T12-C4-B13-D75-NB-007,muc=NB,loigiai={$V=\pi\int_0^{6}\left(\sqrt{x}\right)^2dx=\pi\int_0^{6}1x\,dx=\pi\left[\dfrac{x^2}{2}\right]_0^{6}=\dfrac{36\pi}{2}\approx56{,}55$ m$^3$.}]{Một bồn chứa được tạo thành khi quay quanh trục hoành hình phẳng giới hạn bởi $y=\sqrt{x}$, $y=0$, $x=0$ và $x=6$ (đơn vị trên trục là mét). Thể tích của bồn chứa bằng bao nhiêu mét khối (làm tròn đến hàng phần trăm)?}
  {$56{,}55$.}
  {$113{,}1$.}{$18$.}{$28{,}27$.}

\cauTN[id=T12-C4-B13-D75-NB-008,muc=NB,loigiai={$V=\pi\int_0^{7}\left(\sqrt{4x}\right)^2dx=\pi\int_0^{7}4x\,dx=\pi\left[\dfrac{4x^2}{2}\right]_0^{7}=\dfrac{196\pi}{2}\approx307{,}88$ m$^3$.}]{Một bồn chứa được tạo thành khi quay quanh trục hoành hình phẳng giới hạn bởi $y=\sqrt{4x}$, $y=0$, $x=0$ và $x=7$ (đơn vị trên trục là mét). Thể tích của bồn chứa bằng bao nhiêu mét khối (làm tròn đến hàng phần trăm)?}
  {$307{,}88$.}
  {$615{,}75$.}{$98$.}{$153{,}94$.}

\cauTN[id=T12-C4-B13-D75-NB-009,muc=NB,loigiai={$V=\pi\int_0^{7}\left(\sqrt{2x}\right)^2dx=\pi\int_0^{7}2x\,dx=\pi\left[\dfrac{2x^2}{2}\right]_0^{7}=\dfrac{98\pi}{2}\approx153{,}94$ m$^3$.}]{Một bồn chứa được tạo thành khi quay quanh trục hoành hình phẳng giới hạn bởi $y=\sqrt{2x}$, $y=0$, $x=0$ và $x=7$ (đơn vị trên trục là mét). Thể tích của bồn chứa bằng bao nhiêu mét khối (làm tròn đến hàng phần trăm)?}
  {$153{,}94$.}
  {$307{,}88$.}{$49$.}{$76{,}97$.}

\cauTN[id=T12-C4-B13-D75-NB-010,muc=NB,loigiai={$V=\pi\int_0^{4}\left(\sqrt{x}\right)^2dx=\pi\int_0^{4}1x\,dx=\pi\left[\dfrac{x^2}{2}\right]_0^{4}=\dfrac{16\pi}{2}\approx25{,}13$ m$^3$.}]{Một bồn chứa được tạo thành khi quay quanh trục hoành hình phẳng giới hạn bởi $y=\sqrt{x}$, $y=0$, $x=0$ và $x=4$ (đơn vị trên trục là mét). Thể tích của bồn chứa bằng bao nhiêu mét khối (làm tròn đến hàng phần trăm)?}
  {$25{,}13$.}
  {$50{,}27$.}{$8$.}{$12{,}57$.}

\cauTN[id=T12-C4-B13-D75-NB-011,muc=NB,loigiai={$V=\pi\int_0^{3}\left(\sqrt{2x}\right)^2dx=\pi\int_0^{3}2x\,dx=\pi\left[\dfrac{2x^2}{2}\right]_0^{3}=\dfrac{18\pi}{2}\approx28{,}27$ m$^3$.}]{Một bồn chứa được tạo thành khi quay quanh trục hoành hình phẳng giới hạn bởi $y=\sqrt{2x}$, $y=0$, $x=0$ và $x=3$ (đơn vị trên trục là mét). Thể tích của bồn chứa bằng bao nhiêu mét khối (làm tròn đến hàng phần trăm)?}
  {$28{,}27$.}
  {$56{,}55$.}{$9$.}{$14{,}14$.}

\cauTN[id=T12-C4-B13-D75-NB-012,muc=NB,loigiai={$V=\pi\int_0^{4}\left(\sqrt{4x}\right)^2dx=\pi\int_0^{4}4x\,dx=\pi\left[\dfrac{4x^2}{2}\right]_0^{4}=\dfrac{64\pi}{2}\approx100{,}53$ m$^3$.}]{Một bồn chứa được tạo thành khi quay quanh trục hoành hình phẳng giới hạn bởi $y=\sqrt{4x}$, $y=0$, $x=0$ và $x=4$ (đơn vị trên trục là mét). Thể tích của bồn chứa bằng bao nhiêu mét khối (làm tròn đến hàng phần trăm)?}
  {$100{,}53$.}
  {$201{,}06$.}{$32$.}{$50{,}27$.}

\cauTN[id=T12-C4-B13-D75-NB-013,muc=NB,loigiai={$V=\pi\int_0^{9}\left(\sqrt{3x}\right)^2dx=\pi\int_0^{9}3x\,dx=\pi\left[\dfrac{3x^2}{2}\right]_0^{9}=\dfrac{243\pi}{2}\approx381{,}7$ m$^3$.}]{Một bồn chứa được tạo thành khi quay quanh trục hoành hình phẳng giới hạn bởi $y=\sqrt{3x}$, $y=0$, $x=0$ và $x=9$ (đơn vị trên trục là mét). Thể tích của bồn chứa bằng bao nhiêu mét khối (làm tròn đến hàng phần trăm)?}
  {$381{,}7$.}
  {$763{,}41$.}{$121{,}5$.}{$190{,}85$.}

\cauTN[id=T12-C4-B13-D75-NB-014,muc=NB,loigiai={$V=\pi\int_0^{4}\left(\sqrt{3x}\right)^2dx=\pi\int_0^{4}3x\,dx=\pi\left[\dfrac{3x^2}{2}\right]_0^{4}=\dfrac{48\pi}{2}\approx75{,}4$ m$^3$.}]{Một bồn chứa được tạo thành khi quay quanh trục hoành hình phẳng giới hạn bởi $y=\sqrt{3x}$, $y=0$, $x=0$ và $x=4$ (đơn vị trên trục là mét). Thể tích của bồn chứa bằng bao nhiêu mét khối (làm tròn đến hàng phần trăm)?}
  {$75{,}4$.}
  {$150{,}8$.}{$24$.}{$37{,}7$.}

\cauTN[id=T12-C4-B13-D75-NB-015,muc=NB,loigiai={$V=\pi\int_0^{4}\left(\sqrt{2x}\right)^2dx=\pi\int_0^{4}2x\,dx=\pi\left[\dfrac{2x^2}{2}\right]_0^{4}=\dfrac{32\pi}{2}\approx50{,}27$ m$^3$.}]{Một bồn chứa được tạo thành khi quay quanh trục hoành hình phẳng giới hạn bởi $y=\sqrt{2x}$, $y=0$, $x=0$ và $x=4$ (đơn vị trên trục là mét). Thể tích của bồn chứa bằng bao nhiêu mét khối (làm tròn đến hàng phần trăm)?}
  {$50{,}27$.}
  {$100{,}53$.}{$16$.}{$25{,}13$.}

\cauTN[id=T12-C4-B13-D75-NB-016,muc=NB,loigiai={$V=\pi\int_0^{9}\left(\sqrt{x}\right)^2dx=\pi\int_0^{9}1x\,dx=\pi\left[\dfrac{x^2}{2}\right]_0^{9}=\dfrac{81\pi}{2}\approx127{,}23$ m$^3$.}]{Một bồn chứa được tạo thành khi quay quanh trục hoành hình phẳng giới hạn bởi $y=\sqrt{x}$, $y=0$, $x=0$ và $x=9$ (đơn vị trên trục là mét). Thể tích của bồn chứa bằng bao nhiêu mét khối (làm tròn đến hàng phần trăm)?}
  {$127{,}23$.}
  {$254{,}47$.}{$40{,}5$.}{$63{,}62$.}

\cauTN[id=T12-C4-B13-D75-NB-017,muc=NB,loigiai={$V=\pi\int_0^{5}\left(\sqrt{4x}\right)^2dx=\pi\int_0^{5}4x\,dx=\pi\left[\dfrac{4x^2}{2}\right]_0^{5}=\dfrac{100\pi}{2}\approx157{,}08$ m$^3$.}]{Một bồn chứa được tạo thành khi quay quanh trục hoành hình phẳng giới hạn bởi $y=\sqrt{4x}$, $y=0$, $x=0$ và $x=5$ (đơn vị trên trục là mét). Thể tích của bồn chứa bằng bao nhiêu mét khối (làm tròn đến hàng phần trăm)?}
  {$157{,}08$.}
  {$314{,}16$.}{$50$.}{$78{,}54$.}

\cauTN[id=T12-C4-B13-D75-NB-018,muc=NB,loigiai={$V=\pi\int_0^{3}\left(\sqrt{5x}\right)^2dx=\pi\int_0^{3}5x\,dx=\pi\left[\dfrac{5x^2}{2}\right]_0^{3}=\dfrac{45\pi}{2}\approx70{,}69$ m$^3$.}]{Một bồn chứa được tạo thành khi quay quanh trục hoành hình phẳng giới hạn bởi $y=\sqrt{5x}$, $y=0$, $x=0$ và $x=3$ (đơn vị trên trục là mét). Thể tích của bồn chứa bằng bao nhiêu mét khối (làm tròn đến hàng phần trăm)?}
  {$70{,}69$.}
  {$141{,}37$.}{$22{,}5$.}{$35{,}34$.}

\cauDS[id=T12-C4-B13-D75-NB-019,muc=NB,loigiai={$V=\pi\int_0^{6}\left(\sqrt{3x}\right)^2dx=\pi\int_0^{6}3x\,dx=\pi\left[\dfrac{3x^2}{2}\right]_0^{6}=\dfrac{108\pi}{2}\approx169{,}65$ m$^3$.}]{Một bồn chứa được tạo thành khi quay quanh trục hoành hình phẳng giới hạn bởi $y=\sqrt{3x}$, $y=0$, $x=0$ và $x=6$ (đơn vị trên trục là mét).}{%
  \y{Thể tích được tính bởi $V=\pi\int_0^{6}3x\,dx$.}{d}
  \y{Thể tích bồn bằng $\dfrac{108\pi}{2}$ m$^3$.}{d}
  \y{Thể tích bồn xấp xỉ $169{,}65$ m$^3$.}{d}
  \y{Thể tích bồn bằng $108\pi$ m$^3$.}{s}}

\cauDS[id=T12-C4-B13-D75-NB-020,muc=NB,loigiai={$V=\pi\int_0^{4}\left(\sqrt{3x}\right)^2dx=\pi\int_0^{4}3x\,dx=\pi\left[\dfrac{3x^2}{2}\right]_0^{4}=\dfrac{48\pi}{2}\approx75{,}4$ m$^3$.}]{Một bồn chứa được tạo thành khi quay quanh trục hoành hình phẳng giới hạn bởi $y=\sqrt{3x}$, $y=0$, $x=0$ và $x=4$ (đơn vị trên trục là mét).}{%
  \y{Thể tích được tính bởi $V=\pi\int_0^{4}3x\,dx$.}{d}
  \y{Thể tích bồn bằng $\dfrac{48\pi}{2}$ m$^3$.}{d}
  \y{Thể tích bồn xấp xỉ $75{,}4$ m$^3$.}{d}
  \y{Thể tích bồn bằng $48\pi$ m$^3$.}{s}}

\cauDS[id=T12-C4-B13-D75-NB-021,muc=NB,loigiai={$V=\pi\int_0^{9}\left(\sqrt{5x}\right)^2dx=\pi\int_0^{9}5x\,dx=\pi\left[\dfrac{5x^2}{2}\right]_0^{9}=\dfrac{405\pi}{2}\approx636{,}17$ m$^3$.}]{Một bồn chứa được tạo thành khi quay quanh trục hoành hình phẳng giới hạn bởi $y=\sqrt{5x}$, $y=0$, $x=0$ và $x=9$ (đơn vị trên trục là mét).}{%
  \y{Thể tích được tính bởi $V=\pi\int_0^{9}5x\,dx$.}{d}
  \y{Thể tích bồn bằng $\dfrac{405\pi}{2}$ m$^3$.}{d}
  \y{Thể tích bồn xấp xỉ $636{,}17$ m$^3$.}{d}
  \y{Thể tích bồn bằng $405\pi$ m$^3$.}{s}}

\cauDS[id=T12-C4-B13-D75-NB-022,muc=NB,loigiai={$V=\pi\int_0^{3}\left(\sqrt{x}\right)^2dx=\pi\int_0^{3}1x\,dx=\pi\left[\dfrac{x^2}{2}\right]_0^{3}=\dfrac{9\pi}{2}\approx14{,}14$ m$^3$.}]{Một bồn chứa được tạo thành khi quay quanh trục hoành hình phẳng giới hạn bởi $y=\sqrt{x}$, $y=0$, $x=0$ và $x=3$ (đơn vị trên trục là mét).}{%
  \y{Thể tích được tính bởi $V=\pi\int_0^{3}1x\,dx$.}{d}
  \y{Thể tích bồn bằng $\dfrac{9\pi}{2}$ m$^3$.}{d}
  \y{Thể tích bồn xấp xỉ $14{,}14$ m$^3$.}{d}
  \y{Thể tích bồn bằng $9\pi$ m$^3$.}{s}}

\cauDS[id=T12-C4-B13-D75-NB-023,muc=NB,loigiai={$V=\pi\int_0^{8}\left(\sqrt{4x}\right)^2dx=\pi\int_0^{8}4x\,dx=\pi\left[\dfrac{4x^2}{2}\right]_0^{8}=\dfrac{256\pi}{2}\approx402{,}12$ m$^3$.}]{Một bồn chứa được tạo thành khi quay quanh trục hoành hình phẳng giới hạn bởi $y=\sqrt{4x}$, $y=0$, $x=0$ và $x=8$ (đơn vị trên trục là mét).}{%
  \y{Thể tích được tính bởi $V=\pi\int_0^{8}4x\,dx$.}{d}
  \y{Thể tích bồn bằng $\dfrac{256\pi}{2}$ m$^3$.}{d}
  \y{Thể tích bồn xấp xỉ $402{,}12$ m$^3$.}{d}
  \y{Thể tích bồn bằng $256\pi$ m$^3$.}{s}}

\cauDS[id=T12-C4-B13-D75-NB-024,muc=NB,loigiai={$V=\pi\int_0^{3}\left(\sqrt{3x}\right)^2dx=\pi\int_0^{3}3x\,dx=\pi\left[\dfrac{3x^2}{2}\right]_0^{3}=\dfrac{27\pi}{2}\approx42{,}41$ m$^3$.}]{Một bồn chứa được tạo thành khi quay quanh trục hoành hình phẳng giới hạn bởi $y=\sqrt{3x}$, $y=0$, $x=0$ và $x=3$ (đơn vị trên trục là mét).}{%
  \y{Thể tích được tính bởi $V=\pi\int_0^{3}3x\,dx$.}{d}
  \y{Thể tích bồn bằng $\dfrac{27\pi}{2}$ m$^3$.}{d}
  \y{Thể tích bồn xấp xỉ $42{,}41$ m$^3$.}{d}
  \y{Thể tích bồn bằng $27\pi$ m$^3$.}{s}}

\cauTLN[id=T12-C4-B13-D75-NB-025,muc=NB,loigiai={$V=\pi\int_0^{7}\left(\sqrt{4x}\right)^2dx=\pi\int_0^{7}4x\,dx=\pi\left[\dfrac{4x^2}{2}\right]_0^{7}=\dfrac{196\pi}{2}\approx307{,}88$ m$^3$.}]{Một bồn chứa được tạo thành khi quay quanh trục hoành hình phẳng giới hạn bởi $y=\sqrt{4x}$, $y=0$, $x=0$ và $x=7$ (đơn vị trên trục là mét). Tính thể tích bồn chứa (m$^3$, làm tròn đến hàng phần trăm).}{307,88}

\cauTLN[id=T12-C4-B13-D75-NB-026,muc=NB,loigiai={$V=\pi\int_0^{2}\left(\sqrt{2x}\right)^2dx=\pi\int_0^{2}2x\,dx=\pi\left[\dfrac{2x^2}{2}\right]_0^{2}=\dfrac{8\pi}{2}\approx12{,}57$ m$^3$.}]{Một bồn chứa được tạo thành khi quay quanh trục hoành hình phẳng giới hạn bởi $y=\sqrt{2x}$, $y=0$, $x=0$ và $x=2$ (đơn vị trên trục là mét). Tính thể tích bồn chứa (m$^3$, làm tròn đến hàng phần trăm).}{12,57}

\cauTLN[id=T12-C4-B13-D75-NB-027,muc=NB,loigiai={$V=\pi\int_0^{2}\left(\sqrt{3x}\right)^2dx=\pi\int_0^{2}3x\,dx=\pi\left[\dfrac{3x^2}{2}\right]_0^{2}=\dfrac{12\pi}{2}\approx18{,}85$ m$^3$.}]{Một bồn chứa được tạo thành khi quay quanh trục hoành hình phẳng giới hạn bởi $y=\sqrt{3x}$, $y=0$, $x=0$ và $x=2$ (đơn vị trên trục là mét). Tính thể tích bồn chứa (m$^3$, làm tròn đến hàng phần trăm).}{18,85}

\cauTLN[id=T12-C4-B13-D75-NB-028,muc=NB,loigiai={$V=\pi\int_0^{6}\left(\sqrt{3x}\right)^2dx=\pi\int_0^{6}3x\,dx=\pi\left[\dfrac{3x^2}{2}\right]_0^{6}=\dfrac{108\pi}{2}\approx169{,}65$ m$^3$.}]{Một bồn chứa được tạo thành khi quay quanh trục hoành hình phẳng giới hạn bởi $y=\sqrt{3x}$, $y=0$, $x=0$ và $x=6$ (đơn vị trên trục là mét). Tính thể tích bồn chứa (m$^3$, làm tròn đến hàng phần trăm).}{169,65}

\cauTLN[id=T12-C4-B13-D75-NB-029,muc=NB,loigiai={$V=\pi\int_0^{5}\left(\sqrt{4x}\right)^2dx=\pi\int_0^{5}4x\,dx=\pi\left[\dfrac{4x^2}{2}\right]_0^{5}=\dfrac{100\pi}{2}\approx157{,}08$ m$^3$.}]{Một bồn chứa được tạo thành khi quay quanh trục hoành hình phẳng giới hạn bởi $y=\sqrt{4x}$, $y=0$, $x=0$ và $x=5$ (đơn vị trên trục là mét). Tính thể tích bồn chứa (m$^3$, làm tròn đến hàng phần trăm).}{157,08}

\cauTLN[id=T12-C4-B13-D75-NB-030,muc=NB,loigiai={$V=\pi\int_0^{7}\left(\sqrt{3x}\right)^2dx=\pi\int_0^{7}3x\,dx=\pi\left[\dfrac{3x^2}{2}\right]_0^{7}=\dfrac{147\pi}{2}\approx230{,}91$ m$^3$.}]{Một bồn chứa được tạo thành khi quay quanh trục hoành hình phẳng giới hạn bởi $y=\sqrt{3x}$, $y=0$, $x=0$ và $x=7$ (đơn vị trên trục là mét). Tính thể tích bồn chứa (m$^3$, làm tròn đến hàng phần trăm).}{230,91}

\cauTN[id=T12-C4-B13-D75-TH-001,muc=TH,loigiai={$V=\pi\int_0^{6}\left(\sqrt{5x}\right)^2dx=\pi\int_0^{6}5x\,dx=\pi\left[\dfrac{5x^2}{2}\right]_0^{6}=\dfrac{180\pi}{2}\approx282{,}74$ m$^3$.}]{Một bồn chứa được tạo thành khi quay quanh trục hoành hình phẳng giới hạn bởi $y=\sqrt{5x}$, $y=0$, $x=0$ và $x=6$ (đơn vị trên trục là mét). Thể tích của bồn chứa bằng bao nhiêu mét khối (làm tròn đến hàng phần trăm)?}
  {$282{,}74$.}
  {$565{,}49$.}{$90$.}{$141{,}37$.}

\cauTN[id=T12-C4-B13-D75-TH-002,muc=TH,loigiai={$V=\pi\int_0^{3}\left(\sqrt{3x}\right)^2dx=\pi\int_0^{3}3x\,dx=\pi\left[\dfrac{3x^2}{2}\right]_0^{3}=\dfrac{27\pi}{2}\approx42{,}41$ m$^3$.}]{Một bồn chứa được tạo thành khi quay quanh trục hoành hình phẳng giới hạn bởi $y=\sqrt{3x}$, $y=0$, $x=0$ và $x=3$ (đơn vị trên trục là mét). Thể tích của bồn chứa bằng bao nhiêu mét khối (làm tròn đến hàng phần trăm)?}
  {$42{,}41$.}
  {$84{,}82$.}{$13{,}5$.}{$21{,}21$.}

\cauTN[id=T12-C4-B13-D75-TH-003,muc=TH,loigiai={$V=\pi\int_0^{5}\left(\sqrt{5x}\right)^2dx=\pi\int_0^{5}5x\,dx=\pi\left[\dfrac{5x^2}{2}\right]_0^{5}=\dfrac{125\pi}{2}\approx196{,}35$ m$^3$.}]{Một bồn chứa được tạo thành khi quay quanh trục hoành hình phẳng giới hạn bởi $y=\sqrt{5x}$, $y=0$, $x=0$ và $x=5$ (đơn vị trên trục là mét). Thể tích của bồn chứa bằng bao nhiêu mét khối (làm tròn đến hàng phần trăm)?}
  {$196{,}35$.}
  {$392{,}7$.}{$62{,}5$.}{$98{,}17$.}

\cauTN[id=T12-C4-B13-D75-TH-004,muc=TH,loigiai={$V=\pi\int_0^{7}\left(\sqrt{2x}\right)^2dx=\pi\int_0^{7}2x\,dx=\pi\left[\dfrac{2x^2}{2}\right]_0^{7}=\dfrac{98\pi}{2}\approx153{,}94$ m$^3$.}]{Một bồn chứa được tạo thành khi quay quanh trục hoành hình phẳng giới hạn bởi $y=\sqrt{2x}$, $y=0$, $x=0$ và $x=7$ (đơn vị trên trục là mét). Thể tích của bồn chứa bằng bao nhiêu mét khối (làm tròn đến hàng phần trăm)?}
  {$153{,}94$.}
  {$307{,}88$.}{$49$.}{$76{,}97$.}

\cauTN[id=T12-C4-B13-D75-TH-005,muc=TH,loigiai={$V=\pi\int_0^{8}\left(\sqrt{3x}\right)^2dx=\pi\int_0^{8}3x\,dx=\pi\left[\dfrac{3x^2}{2}\right]_0^{8}=\dfrac{192\pi}{2}\approx301{,}59$ m$^3$.}]{Một bồn chứa được tạo thành khi quay quanh trục hoành hình phẳng giới hạn bởi $y=\sqrt{3x}$, $y=0$, $x=0$ và $x=8$ (đơn vị trên trục là mét). Thể tích của bồn chứa bằng bao nhiêu mét khối (làm tròn đến hàng phần trăm)?}
  {$301{,}59$.}
  {$603{,}19$.}{$96$.}{$150{,}8$.}

\cauTN[id=T12-C4-B13-D75-TH-006,muc=TH,loigiai={$V=\pi\int_0^{4}\left(\sqrt{x}\right)^2dx=\pi\int_0^{4}1x\,dx=\pi\left[\dfrac{x^2}{2}\right]_0^{4}=\dfrac{16\pi}{2}\approx25{,}13$ m$^3$.}]{Một bồn chứa được tạo thành khi quay quanh trục hoành hình phẳng giới hạn bởi $y=\sqrt{x}$, $y=0$, $x=0$ và $x=4$ (đơn vị trên trục là mét). Thể tích của bồn chứa bằng bao nhiêu mét khối (làm tròn đến hàng phần trăm)?}
  {$25{,}13$.}
  {$50{,}27$.}{$8$.}{$12{,}57$.}

\cauTN[id=T12-C4-B13-D75-TH-007,muc=TH,loigiai={$V=\pi\int_0^{9}\left(\sqrt{2x}\right)^2dx=\pi\int_0^{9}2x\,dx=\pi\left[\dfrac{2x^2}{2}\right]_0^{9}=\dfrac{162\pi}{2}\approx254{,}47$ m$^3$.}]{Một bồn chứa được tạo thành khi quay quanh trục hoành hình phẳng giới hạn bởi $y=\sqrt{2x}$, $y=0$, $x=0$ và $x=9$ (đơn vị trên trục là mét). Thể tích của bồn chứa bằng bao nhiêu mét khối (làm tròn đến hàng phần trăm)?}
  {$254{,}47$.}
  {$508{,}94$.}{$81$.}{$127{,}23$.}

\cauTN[id=T12-C4-B13-D75-TH-008,muc=TH,loigiai={$V=\pi\int_0^{8}\left(\sqrt{5x}\right)^2dx=\pi\int_0^{8}5x\,dx=\pi\left[\dfrac{5x^2}{2}\right]_0^{8}=\dfrac{320\pi}{2}\approx502{,}65$ m$^3$.}]{Một bồn chứa được tạo thành khi quay quanh trục hoành hình phẳng giới hạn bởi $y=\sqrt{5x}$, $y=0$, $x=0$ và $x=8$ (đơn vị trên trục là mét). Thể tích của bồn chứa bằng bao nhiêu mét khối (làm tròn đến hàng phần trăm)?}
  {$502{,}65$.}
  {$1005{,}31$.}{$160$.}{$251{,}33$.}

\cauTN[id=T12-C4-B13-D75-TH-009,muc=TH,loigiai={$V=\pi\int_0^{7}\left(\sqrt{2x}\right)^2dx=\pi\int_0^{7}2x\,dx=\pi\left[\dfrac{2x^2}{2}\right]_0^{7}=\dfrac{98\pi}{2}\approx153{,}94$ m$^3$.}]{Một bồn chứa được tạo thành khi quay quanh trục hoành hình phẳng giới hạn bởi $y=\sqrt{2x}$, $y=0$, $x=0$ và $x=7$ (đơn vị trên trục là mét). Thể tích của bồn chứa bằng bao nhiêu mét khối (làm tròn đến hàng phần trăm)?}
  {$153{,}94$.}
  {$307{,}88$.}{$49$.}{$76{,}97$.}

\cauTN[id=T12-C4-B13-D75-TH-010,muc=TH,loigiai={$V=\pi\int_0^{6}\left(\sqrt{2x}\right)^2dx=\pi\int_0^{6}2x\,dx=\pi\left[\dfrac{2x^2}{2}\right]_0^{6}=\dfrac{72\pi}{2}\approx113{,}1$ m$^3$.}]{Một bồn chứa được tạo thành khi quay quanh trục hoành hình phẳng giới hạn bởi $y=\sqrt{2x}$, $y=0$, $x=0$ và $x=6$ (đơn vị trên trục là mét). Thể tích của bồn chứa bằng bao nhiêu mét khối (làm tròn đến hàng phần trăm)?}
  {$113{,}1$.}
  {$226{,}19$.}{$36$.}{$56{,}55$.}

\cauTN[id=T12-C4-B13-D75-TH-011,muc=TH,loigiai={$V=\pi\int_0^{8}\left(\sqrt{4x}\right)^2dx=\pi\int_0^{8}4x\,dx=\pi\left[\dfrac{4x^2}{2}\right]_0^{8}=\dfrac{256\pi}{2}\approx402{,}12$ m$^3$.}]{Một bồn chứa được tạo thành khi quay quanh trục hoành hình phẳng giới hạn bởi $y=\sqrt{4x}$, $y=0$, $x=0$ và $x=8$ (đơn vị trên trục là mét). Thể tích của bồn chứa bằng bao nhiêu mét khối (làm tròn đến hàng phần trăm)?}
  {$402{,}12$.}
  {$804{,}25$.}{$128$.}{$201{,}06$.}

\cauTN[id=T12-C4-B13-D75-TH-012,muc=TH,loigiai={$V=\pi\int_0^{4}\left(\sqrt{x}\right)^2dx=\pi\int_0^{4}1x\,dx=\pi\left[\dfrac{x^2}{2}\right]_0^{4}=\dfrac{16\pi}{2}\approx25{,}13$ m$^3$.}]{Một bồn chứa được tạo thành khi quay quanh trục hoành hình phẳng giới hạn bởi $y=\sqrt{x}$, $y=0$, $x=0$ và $x=4$ (đơn vị trên trục là mét). Thể tích của bồn chứa bằng bao nhiêu mét khối (làm tròn đến hàng phần trăm)?}
  {$25{,}13$.}
  {$50{,}27$.}{$8$.}{$12{,}57$.}

\cauTN[id=T12-C4-B13-D75-TH-013,muc=TH,loigiai={$V=\pi\int_0^{6}\left(\sqrt{3x}\right)^2dx=\pi\int_0^{6}3x\,dx=\pi\left[\dfrac{3x^2}{2}\right]_0^{6}=\dfrac{108\pi}{2}\approx169{,}65$ m$^3$.}]{Một bồn chứa được tạo thành khi quay quanh trục hoành hình phẳng giới hạn bởi $y=\sqrt{3x}$, $y=0$, $x=0$ và $x=6$ (đơn vị trên trục là mét). Thể tích của bồn chứa bằng bao nhiêu mét khối (làm tròn đến hàng phần trăm)?}
  {$169{,}65$.}
  {$339{,}29$.}{$54$.}{$84{,}82$.}

\cauTN[id=T12-C4-B13-D75-TH-014,muc=TH,loigiai={$V=\pi\int_0^{5}\left(\sqrt{3x}\right)^2dx=\pi\int_0^{5}3x\,dx=\pi\left[\dfrac{3x^2}{2}\right]_0^{5}=\dfrac{75\pi}{2}\approx117{,}81$ m$^3$.}]{Một bồn chứa được tạo thành khi quay quanh trục hoành hình phẳng giới hạn bởi $y=\sqrt{3x}$, $y=0$, $x=0$ và $x=5$ (đơn vị trên trục là mét). Thể tích của bồn chứa bằng bao nhiêu mét khối (làm tròn đến hàng phần trăm)?}
  {$117{,}81$.}
  {$235{,}62$.}{$37{,}5$.}{$58{,}9$.}

\cauTN[id=T12-C4-B13-D75-TH-015,muc=TH,loigiai={$V=\pi\int_0^{3}\left(\sqrt{2x}\right)^2dx=\pi\int_0^{3}2x\,dx=\pi\left[\dfrac{2x^2}{2}\right]_0^{3}=\dfrac{18\pi}{2}\approx28{,}27$ m$^3$.}]{Một bồn chứa được tạo thành khi quay quanh trục hoành hình phẳng giới hạn bởi $y=\sqrt{2x}$, $y=0$, $x=0$ và $x=3$ (đơn vị trên trục là mét). Thể tích của bồn chứa bằng bao nhiêu mét khối (làm tròn đến hàng phần trăm)?}
  {$28{,}27$.}
  {$56{,}55$.}{$9$.}{$14{,}14$.}

\cauTN[id=T12-C4-B13-D75-TH-016,muc=TH,loigiai={$V=\pi\int_0^{4}\left(\sqrt{2x}\right)^2dx=\pi\int_0^{4}2x\,dx=\pi\left[\dfrac{2x^2}{2}\right]_0^{4}=\dfrac{32\pi}{2}\approx50{,}27$ m$^3$.}]{Một bồn chứa được tạo thành khi quay quanh trục hoành hình phẳng giới hạn bởi $y=\sqrt{2x}$, $y=0$, $x=0$ và $x=4$ (đơn vị trên trục là mét). Thể tích của bồn chứa bằng bao nhiêu mét khối (làm tròn đến hàng phần trăm)?}
  {$50{,}27$.}
  {$100{,}53$.}{$16$.}{$25{,}13$.}

\cauTN[id=T12-C4-B13-D75-TH-017,muc=TH,loigiai={$V=\pi\int_0^{5}\left(\sqrt{5x}\right)^2dx=\pi\int_0^{5}5x\,dx=\pi\left[\dfrac{5x^2}{2}\right]_0^{5}=\dfrac{125\pi}{2}\approx196{,}35$ m$^3$.}]{Một bồn chứa được tạo thành khi quay quanh trục hoành hình phẳng giới hạn bởi $y=\sqrt{5x}$, $y=0$, $x=0$ và $x=5$ (đơn vị trên trục là mét). Thể tích của bồn chứa bằng bao nhiêu mét khối (làm tròn đến hàng phần trăm)?}
  {$196{,}35$.}
  {$392{,}7$.}{$62{,}5$.}{$98{,}17$.}

\cauTN[id=T12-C4-B13-D75-TH-018,muc=TH,loigiai={$V=\pi\int_0^{9}\left(\sqrt{5x}\right)^2dx=\pi\int_0^{9}5x\,dx=\pi\left[\dfrac{5x^2}{2}\right]_0^{9}=\dfrac{405\pi}{2}\approx636{,}17$ m$^3$.}]{Một bồn chứa được tạo thành khi quay quanh trục hoành hình phẳng giới hạn bởi $y=\sqrt{5x}$, $y=0$, $x=0$ và $x=9$ (đơn vị trên trục là mét). Thể tích của bồn chứa bằng bao nhiêu mét khối (làm tròn đến hàng phần trăm)?}
  {$636{,}17$.}
  {$1272{,}35$.}{$202{,}5$.}{$318{,}09$.}

\cauTN[id=T12-C4-B13-D75-TH-019,muc=TH,loigiai={$V=\pi\int_0^{9}\left(\sqrt{4x}\right)^2dx=\pi\int_0^{9}4x\,dx=\pi\left[\dfrac{4x^2}{2}\right]_0^{9}=\dfrac{324\pi}{2}\approx508{,}94$ m$^3$.}]{Một bồn chứa được tạo thành khi quay quanh trục hoành hình phẳng giới hạn bởi $y=\sqrt{4x}$, $y=0$, $x=0$ và $x=9$ (đơn vị trên trục là mét). Thể tích của bồn chứa bằng bao nhiêu mét khối (làm tròn đến hàng phần trăm)?}
  {$508{,}94$.}
  {$1017{,}88$.}{$162$.}{$254{,}47$.}

\cauTN[id=T12-C4-B13-D75-TH-020,muc=TH,loigiai={$V=\pi\int_0^{6}\left(\sqrt{5x}\right)^2dx=\pi\int_0^{6}5x\,dx=\pi\left[\dfrac{5x^2}{2}\right]_0^{6}=\dfrac{180\pi}{2}\approx282{,}74$ m$^3$.}]{Một bồn chứa được tạo thành khi quay quanh trục hoành hình phẳng giới hạn bởi $y=\sqrt{5x}$, $y=0$, $x=0$ và $x=6$ (đơn vị trên trục là mét). Thể tích của bồn chứa bằng bao nhiêu mét khối (làm tròn đến hàng phần trăm)?}
  {$282{,}74$.}
  {$565{,}49$.}{$90$.}{$141{,}37$.}

\cauTN[id=T12-C4-B13-D75-TH-021,muc=TH,loigiai={$V=\pi\int_0^{8}\left(\sqrt{x}\right)^2dx=\pi\int_0^{8}1x\,dx=\pi\left[\dfrac{x^2}{2}\right]_0^{8}=\dfrac{64\pi}{2}\approx100{,}53$ m$^3$.}]{Một bồn chứa được tạo thành khi quay quanh trục hoành hình phẳng giới hạn bởi $y=\sqrt{x}$, $y=0$, $x=0$ và $x=8$ (đơn vị trên trục là mét). Thể tích của bồn chứa bằng bao nhiêu mét khối (làm tròn đến hàng phần trăm)?}
  {$100{,}53$.}
  {$201{,}06$.}{$32$.}{$50{,}27$.}

\cauTN[id=T12-C4-B13-D75-TH-022,muc=TH,loigiai={$V=\pi\int_0^{2}\left(\sqrt{3x}\right)^2dx=\pi\int_0^{2}3x\,dx=\pi\left[\dfrac{3x^2}{2}\right]_0^{2}=\dfrac{12\pi}{2}\approx18{,}85$ m$^3$.}]{Một bồn chứa được tạo thành khi quay quanh trục hoành hình phẳng giới hạn bởi $y=\sqrt{3x}$, $y=0$, $x=0$ và $x=2$ (đơn vị trên trục là mét). Thể tích của bồn chứa bằng bao nhiêu mét khối (làm tròn đến hàng phần trăm)?}
  {$18{,}85$.}
  {$37{,}7$.}{$6$.}{$9{,}42$.}

\cauTN[id=T12-C4-B13-D75-TH-023,muc=TH,loigiai={$V=\pi\int_0^{5}\left(\sqrt{4x}\right)^2dx=\pi\int_0^{5}4x\,dx=\pi\left[\dfrac{4x^2}{2}\right]_0^{5}=\dfrac{100\pi}{2}\approx157{,}08$ m$^3$.}]{Một bồn chứa được tạo thành khi quay quanh trục hoành hình phẳng giới hạn bởi $y=\sqrt{4x}$, $y=0$, $x=0$ và $x=5$ (đơn vị trên trục là mét). Thể tích của bồn chứa bằng bao nhiêu mét khối (làm tròn đến hàng phần trăm)?}
  {$157{,}08$.}
  {$314{,}16$.}{$50$.}{$78{,}54$.}

\cauTN[id=T12-C4-B13-D75-TH-024,muc=TH,loigiai={$V=\pi\int_0^{4}\left(\sqrt{3x}\right)^2dx=\pi\int_0^{4}3x\,dx=\pi\left[\dfrac{3x^2}{2}\right]_0^{4}=\dfrac{48\pi}{2}\approx75{,}4$ m$^3$.}]{Một bồn chứa được tạo thành khi quay quanh trục hoành hình phẳng giới hạn bởi $y=\sqrt{3x}$, $y=0$, $x=0$ và $x=4$ (đơn vị trên trục là mét). Thể tích của bồn chứa bằng bao nhiêu mét khối (làm tròn đến hàng phần trăm)?}
  {$75{,}4$.}
  {$150{,}8$.}{$24$.}{$37{,}7$.}

\cauDS[id=T12-C4-B13-D75-TH-025,muc=TH,loigiai={$V=\pi\int_0^{6}\left(\sqrt{x}\right)^2dx=\pi\int_0^{6}1x\,dx=\pi\left[\dfrac{x^2}{2}\right]_0^{6}=\dfrac{36\pi}{2}\approx56{,}55$ m$^3$.}]{Một bồn chứa được tạo thành khi quay quanh trục hoành hình phẳng giới hạn bởi $y=\sqrt{x}$, $y=0$, $x=0$ và $x=6$ (đơn vị trên trục là mét).}{%
  \y{Thể tích được tính bởi $V=\pi\int_0^{6}1x\,dx$.}{d}
  \y{Thể tích bồn bằng $\dfrac{36\pi}{2}$ m$^3$.}{d}
  \y{Thể tích bồn xấp xỉ $56{,}55$ m$^3$.}{d}
  \y{Thể tích bồn bằng $36\pi$ m$^3$.}{s}}

\cauDS[id=T12-C4-B13-D75-TH-026,muc=TH,loigiai={$V=\pi\int_0^{7}\left(\sqrt{5x}\right)^2dx=\pi\int_0^{7}5x\,dx=\pi\left[\dfrac{5x^2}{2}\right]_0^{7}=\dfrac{245\pi}{2}\approx384{,}85$ m$^3$.}]{Một bồn chứa được tạo thành khi quay quanh trục hoành hình phẳng giới hạn bởi $y=\sqrt{5x}$, $y=0$, $x=0$ và $x=7$ (đơn vị trên trục là mét).}{%
  \y{Thể tích được tính bởi $V=\pi\int_0^{7}5x\,dx$.}{d}
  \y{Thể tích bồn bằng $\dfrac{245\pi}{2}$ m$^3$.}{d}
  \y{Thể tích bồn xấp xỉ $384{,}85$ m$^3$.}{d}
  \y{Thể tích bồn bằng $245\pi$ m$^3$.}{s}}

\cauDS[id=T12-C4-B13-D75-TH-027,muc=TH,loigiai={$V=\pi\int_0^{6}\left(\sqrt{3x}\right)^2dx=\pi\int_0^{6}3x\,dx=\pi\left[\dfrac{3x^2}{2}\right]_0^{6}=\dfrac{108\pi}{2}\approx169{,}65$ m$^3$.}]{Một bồn chứa được tạo thành khi quay quanh trục hoành hình phẳng giới hạn bởi $y=\sqrt{3x}$, $y=0$, $x=0$ và $x=6$ (đơn vị trên trục là mét).}{%
  \y{Thể tích được tính bởi $V=\pi\int_0^{6}3x\,dx$.}{d}
  \y{Thể tích bồn bằng $\dfrac{108\pi}{2}$ m$^3$.}{d}
  \y{Thể tích bồn xấp xỉ $169{,}65$ m$^3$.}{d}
  \y{Thể tích bồn bằng $108\pi$ m$^3$.}{s}}

\cauDS[id=T12-C4-B13-D75-TH-028,muc=TH,loigiai={$V=\pi\int_0^{6}\left(\sqrt{4x}\right)^2dx=\pi\int_0^{6}4x\,dx=\pi\left[\dfrac{4x^2}{2}\right]_0^{6}=\dfrac{144\pi}{2}\approx226{,}19$ m$^3$.}]{Một bồn chứa được tạo thành khi quay quanh trục hoành hình phẳng giới hạn bởi $y=\sqrt{4x}$, $y=0$, $x=0$ và $x=6$ (đơn vị trên trục là mét).}{%
  \y{Thể tích được tính bởi $V=\pi\int_0^{6}4x\,dx$.}{d}
  \y{Thể tích bồn bằng $\dfrac{144\pi}{2}$ m$^3$.}{d}
  \y{Thể tích bồn xấp xỉ $226{,}19$ m$^3$.}{d}
  \y{Thể tích bồn bằng $144\pi$ m$^3$.}{s}}

\cauDS[id=T12-C4-B13-D75-TH-029,muc=TH,loigiai={$V=\pi\int_0^{9}\left(\sqrt{2x}\right)^2dx=\pi\int_0^{9}2x\,dx=\pi\left[\dfrac{2x^2}{2}\right]_0^{9}=\dfrac{162\pi}{2}\approx254{,}47$ m$^3$.}]{Một bồn chứa được tạo thành khi quay quanh trục hoành hình phẳng giới hạn bởi $y=\sqrt{2x}$, $y=0$, $x=0$ và $x=9$ (đơn vị trên trục là mét).}{%
  \y{Thể tích được tính bởi $V=\pi\int_0^{9}2x\,dx$.}{d}
  \y{Thể tích bồn bằng $\dfrac{162\pi}{2}$ m$^3$.}{d}
  \y{Thể tích bồn xấp xỉ $254{,}47$ m$^3$.}{d}
  \y{Thể tích bồn bằng $162\pi$ m$^3$.}{s}}

\cauDS[id=T12-C4-B13-D75-TH-030,muc=TH,loigiai={$V=\pi\int_0^{3}\left(\sqrt{2x}\right)^2dx=\pi\int_0^{3}2x\,dx=\pi\left[\dfrac{2x^2}{2}\right]_0^{3}=\dfrac{18\pi}{2}\approx28{,}27$ m$^3$.}]{Một bồn chứa được tạo thành khi quay quanh trục hoành hình phẳng giới hạn bởi $y=\sqrt{2x}$, $y=0$, $x=0$ và $x=3$ (đơn vị trên trục là mét).}{%
  \y{Thể tích được tính bởi $V=\pi\int_0^{3}2x\,dx$.}{d}
  \y{Thể tích bồn bằng $\dfrac{18\pi}{2}$ m$^3$.}{d}
  \y{Thể tích bồn xấp xỉ $28{,}27$ m$^3$.}{d}
  \y{Thể tích bồn bằng $18\pi$ m$^3$.}{s}}

\cauDS[id=T12-C4-B13-D75-TH-031,muc=TH,loigiai={$V=\pi\int_0^{5}\left(\sqrt{2x}\right)^2dx=\pi\int_0^{5}2x\,dx=\pi\left[\dfrac{2x^2}{2}\right]_0^{5}=\dfrac{50\pi}{2}\approx78{,}54$ m$^3$.}]{Một bồn chứa được tạo thành khi quay quanh trục hoành hình phẳng giới hạn bởi $y=\sqrt{2x}$, $y=0$, $x=0$ và $x=5$ (đơn vị trên trục là mét).}{%
  \y{Thể tích được tính bởi $V=\pi\int_0^{5}2x\,dx$.}{d}
  \y{Thể tích bồn bằng $\dfrac{50\pi}{2}$ m$^3$.}{d}
  \y{Thể tích bồn xấp xỉ $78{,}54$ m$^3$.}{d}
  \y{Thể tích bồn bằng $50\pi$ m$^3$.}{s}}

\cauDS[id=T12-C4-B13-D75-TH-032,muc=TH,loigiai={$V=\pi\int_0^{3}\left(\sqrt{3x}\right)^2dx=\pi\int_0^{3}3x\,dx=\pi\left[\dfrac{3x^2}{2}\right]_0^{3}=\dfrac{27\pi}{2}\approx42{,}41$ m$^3$.}]{Một bồn chứa được tạo thành khi quay quanh trục hoành hình phẳng giới hạn bởi $y=\sqrt{3x}$, $y=0$, $x=0$ và $x=3$ (đơn vị trên trục là mét).}{%
  \y{Thể tích được tính bởi $V=\pi\int_0^{3}3x\,dx$.}{d}
  \y{Thể tích bồn bằng $\dfrac{27\pi}{2}$ m$^3$.}{d}
  \y{Thể tích bồn xấp xỉ $42{,}41$ m$^3$.}{d}
  \y{Thể tích bồn bằng $27\pi$ m$^3$.}{s}}

\cauTLN[id=T12-C4-B13-D75-TH-033,muc=TH,loigiai={$V=\pi\int_0^{8}\left(\sqrt{4x}\right)^2dx=\pi\int_0^{8}4x\,dx=\pi\left[\dfrac{4x^2}{2}\right]_0^{8}=\dfrac{256\pi}{2}\approx402{,}12$ m$^3$.}]{Một bồn chứa được tạo thành khi quay quanh trục hoành hình phẳng giới hạn bởi $y=\sqrt{4x}$, $y=0$, $x=0$ và $x=8$ (đơn vị trên trục là mét). Tính thể tích bồn chứa (m$^3$, làm tròn đến hàng phần trăm).}{402,12}

\cauTLN[id=T12-C4-B13-D75-TH-034,muc=TH,loigiai={$V=\pi\int_0^{2}\left(\sqrt{2x}\right)^2dx=\pi\int_0^{2}2x\,dx=\pi\left[\dfrac{2x^2}{2}\right]_0^{2}=\dfrac{8\pi}{2}\approx12{,}57$ m$^3$.}]{Một bồn chứa được tạo thành khi quay quanh trục hoành hình phẳng giới hạn bởi $y=\sqrt{2x}$, $y=0$, $x=0$ và $x=2$ (đơn vị trên trục là mét). Tính thể tích bồn chứa (m$^3$, làm tròn đến hàng phần trăm).}{12,57}

\cauTLN[id=T12-C4-B13-D75-TH-035,muc=TH,loigiai={$V=\pi\int_0^{3}\left(\sqrt{4x}\right)^2dx=\pi\int_0^{3}4x\,dx=\pi\left[\dfrac{4x^2}{2}\right]_0^{3}=\dfrac{36\pi}{2}\approx56{,}55$ m$^3$.}]{Một bồn chứa được tạo thành khi quay quanh trục hoành hình phẳng giới hạn bởi $y=\sqrt{4x}$, $y=0$, $x=0$ và $x=3$ (đơn vị trên trục là mét). Tính thể tích bồn chứa (m$^3$, làm tròn đến hàng phần trăm).}{56,55}

\cauTLN[id=T12-C4-B13-D75-TH-036,muc=TH,loigiai={$V=\pi\int_0^{3}\left(\sqrt{2x}\right)^2dx=\pi\int_0^{3}2x\,dx=\pi\left[\dfrac{2x^2}{2}\right]_0^{3}=\dfrac{18\pi}{2}\approx28{,}27$ m$^3$.}]{Một bồn chứa được tạo thành khi quay quanh trục hoành hình phẳng giới hạn bởi $y=\sqrt{2x}$, $y=0$, $x=0$ và $x=3$ (đơn vị trên trục là mét). Tính thể tích bồn chứa (m$^3$, làm tròn đến hàng phần trăm).}{28,27}

\cauTLN[id=T12-C4-B13-D75-TH-037,muc=TH,loigiai={$V=\pi\int_0^{8}\left(\sqrt{4x}\right)^2dx=\pi\int_0^{8}4x\,dx=\pi\left[\dfrac{4x^2}{2}\right]_0^{8}=\dfrac{256\pi}{2}\approx402{,}12$ m$^3$.}]{Một bồn chứa được tạo thành khi quay quanh trục hoành hình phẳng giới hạn bởi $y=\sqrt{4x}$, $y=0$, $x=0$ và $x=8$ (đơn vị trên trục là mét). Tính thể tích bồn chứa (m$^3$, làm tròn đến hàng phần trăm).}{402,12}

\cauTLN[id=T12-C4-B13-D75-TH-038,muc=TH,loigiai={$V=\pi\int_0^{4}\left(\sqrt{3x}\right)^2dx=\pi\int_0^{4}3x\,dx=\pi\left[\dfrac{3x^2}{2}\right]_0^{4}=\dfrac{48\pi}{2}\approx75{,}4$ m$^3$.}]{Một bồn chứa được tạo thành khi quay quanh trục hoành hình phẳng giới hạn bởi $y=\sqrt{3x}$, $y=0$, $x=0$ và $x=4$ (đơn vị trên trục là mét). Tính thể tích bồn chứa (m$^3$, làm tròn đến hàng phần trăm).}{75,4}

\cauTLN[id=T12-C4-B13-D75-TH-039,muc=TH,loigiai={$V=\pi\int_0^{8}\left(\sqrt{2x}\right)^2dx=\pi\int_0^{8}2x\,dx=\pi\left[\dfrac{2x^2}{2}\right]_0^{8}=\dfrac{128\pi}{2}\approx201{,}06$ m$^3$.}]{Một bồn chứa được tạo thành khi quay quanh trục hoành hình phẳng giới hạn bởi $y=\sqrt{2x}$, $y=0$, $x=0$ và $x=8$ (đơn vị trên trục là mét). Tính thể tích bồn chứa (m$^3$, làm tròn đến hàng phần trăm).}{201,06}

\cauTLN[id=T12-C4-B13-D75-TH-040,muc=TH,loigiai={$V=\pi\int_0^{9}\left(\sqrt{5x}\right)^2dx=\pi\int_0^{9}5x\,dx=\pi\left[\dfrac{5x^2}{2}\right]_0^{9}=\dfrac{405\pi}{2}\approx636{,}17$ m$^3$.}]{Một bồn chứa được tạo thành khi quay quanh trục hoành hình phẳng giới hạn bởi $y=\sqrt{5x}$, $y=0$, $x=0$ và $x=9$ (đơn vị trên trục là mét). Tính thể tích bồn chứa (m$^3$, làm tròn đến hàng phần trăm).}{636,17}

\cauTN[id=T12-C4-B13-D75-VD-001,muc=VD,loigiai={$V=\pi\int_0^{6}\left(\sqrt{4x}\right)^2dx=\pi\int_0^{6}4x\,dx=\pi\left[\dfrac{4x^2}{2}\right]_0^{6}=\dfrac{144\pi}{2}\approx226{,}19$ m$^3$.}]{Một bồn chứa được tạo thành khi quay quanh trục hoành hình phẳng giới hạn bởi $y=\sqrt{4x}$, $y=0$, $x=0$ và $x=6$ (đơn vị trên trục là mét). Thể tích của bồn chứa bằng bao nhiêu mét khối (làm tròn đến hàng phần trăm)?}
  {$226{,}19$.}
  {$452{,}39$.}{$72$.}{$113{,}1$.}

\cauTN[id=T12-C4-B13-D75-VD-002,muc=VD,loigiai={$V=\pi\int_0^{5}\left(\sqrt{5x}\right)^2dx=\pi\int_0^{5}5x\,dx=\pi\left[\dfrac{5x^2}{2}\right]_0^{5}=\dfrac{125\pi}{2}\approx196{,}35$ m$^3$.}]{Một bồn chứa được tạo thành khi quay quanh trục hoành hình phẳng giới hạn bởi $y=\sqrt{5x}$, $y=0$, $x=0$ và $x=5$ (đơn vị trên trục là mét). Thể tích của bồn chứa bằng bao nhiêu mét khối (làm tròn đến hàng phần trăm)?}
  {$196{,}35$.}
  {$392{,}7$.}{$62{,}5$.}{$98{,}17$.}

\cauTN[id=T12-C4-B13-D75-VD-003,muc=VD,loigiai={$V=\pi\int_0^{2}\left(\sqrt{x}\right)^2dx=\pi\int_0^{2}1x\,dx=\pi\left[\dfrac{x^2}{2}\right]_0^{2}=\dfrac{4\pi}{2}\approx6{,}28$ m$^3$.}]{Một bồn chứa được tạo thành khi quay quanh trục hoành hình phẳng giới hạn bởi $y=\sqrt{x}$, $y=0$, $x=0$ và $x=2$ (đơn vị trên trục là mét). Thể tích của bồn chứa bằng bao nhiêu mét khối (làm tròn đến hàng phần trăm)?}
  {$6{,}28$.}
  {$12{,}57$.}{$2$.}{$3{,}14$.}

\cauTN[id=T12-C4-B13-D75-VD-004,muc=VD,loigiai={$V=\pi\int_0^{4}\left(\sqrt{3x}\right)^2dx=\pi\int_0^{4}3x\,dx=\pi\left[\dfrac{3x^2}{2}\right]_0^{4}=\dfrac{48\pi}{2}\approx75{,}4$ m$^3$.}]{Một bồn chứa được tạo thành khi quay quanh trục hoành hình phẳng giới hạn bởi $y=\sqrt{3x}$, $y=0$, $x=0$ và $x=4$ (đơn vị trên trục là mét). Thể tích của bồn chứa bằng bao nhiêu mét khối (làm tròn đến hàng phần trăm)?}
  {$75{,}4$.}
  {$150{,}8$.}{$24$.}{$37{,}7$.}

\cauTN[id=T12-C4-B13-D75-VD-005,muc=VD,loigiai={$V=\pi\int_0^{2}\left(\sqrt{2x}\right)^2dx=\pi\int_0^{2}2x\,dx=\pi\left[\dfrac{2x^2}{2}\right]_0^{2}=\dfrac{8\pi}{2}\approx12{,}57$ m$^3$.}]{Một bồn chứa được tạo thành khi quay quanh trục hoành hình phẳng giới hạn bởi $y=\sqrt{2x}$, $y=0$, $x=0$ và $x=2$ (đơn vị trên trục là mét). Thể tích của bồn chứa bằng bao nhiêu mét khối (làm tròn đến hàng phần trăm)?}
  {$12{,}57$.}
  {$25{,}13$.}{$4$.}{$6{,}28$.}

\cauTN[id=T12-C4-B13-D75-VD-006,muc=VD,loigiai={$V=\pi\int_0^{4}\left(\sqrt{3x}\right)^2dx=\pi\int_0^{4}3x\,dx=\pi\left[\dfrac{3x^2}{2}\right]_0^{4}=\dfrac{48\pi}{2}\approx75{,}4$ m$^3$.}]{Một bồn chứa được tạo thành khi quay quanh trục hoành hình phẳng giới hạn bởi $y=\sqrt{3x}$, $y=0$, $x=0$ và $x=4$ (đơn vị trên trục là mét). Thể tích của bồn chứa bằng bao nhiêu mét khối (làm tròn đến hàng phần trăm)?}
  {$75{,}4$.}
  {$150{,}8$.}{$24$.}{$37{,}7$.}

\cauTN[id=T12-C4-B13-D75-VD-007,muc=VD,loigiai={$V=\pi\int_0^{8}\left(\sqrt{x}\right)^2dx=\pi\int_0^{8}1x\,dx=\pi\left[\dfrac{x^2}{2}\right]_0^{8}=\dfrac{64\pi}{2}\approx100{,}53$ m$^3$.}]{Một bồn chứa được tạo thành khi quay quanh trục hoành hình phẳng giới hạn bởi $y=\sqrt{x}$, $y=0$, $x=0$ và $x=8$ (đơn vị trên trục là mét). Thể tích của bồn chứa bằng bao nhiêu mét khối (làm tròn đến hàng phần trăm)?}
  {$100{,}53$.}
  {$201{,}06$.}{$32$.}{$50{,}27$.}

\cauTN[id=T12-C4-B13-D75-VD-008,muc=VD,loigiai={$V=\pi\int_0^{6}\left(\sqrt{x}\right)^2dx=\pi\int_0^{6}1x\,dx=\pi\left[\dfrac{x^2}{2}\right]_0^{6}=\dfrac{36\pi}{2}\approx56{,}55$ m$^3$.}]{Một bồn chứa được tạo thành khi quay quanh trục hoành hình phẳng giới hạn bởi $y=\sqrt{x}$, $y=0$, $x=0$ và $x=6$ (đơn vị trên trục là mét). Thể tích của bồn chứa bằng bao nhiêu mét khối (làm tròn đến hàng phần trăm)?}
  {$56{,}55$.}
  {$113{,}1$.}{$18$.}{$28{,}27$.}

\cauTN[id=T12-C4-B13-D75-VD-009,muc=VD,loigiai={$V=\pi\int_0^{3}\left(\sqrt{4x}\right)^2dx=\pi\int_0^{3}4x\,dx=\pi\left[\dfrac{4x^2}{2}\right]_0^{3}=\dfrac{36\pi}{2}\approx56{,}55$ m$^3$.}]{Một bồn chứa được tạo thành khi quay quanh trục hoành hình phẳng giới hạn bởi $y=\sqrt{4x}$, $y=0$, $x=0$ và $x=3$ (đơn vị trên trục là mét). Thể tích của bồn chứa bằng bao nhiêu mét khối (làm tròn đến hàng phần trăm)?}
  {$56{,}55$.}
  {$113{,}1$.}{$18$.}{$28{,}27$.}

\cauTN[id=T12-C4-B13-D75-VD-010,muc=VD,loigiai={$V=\pi\int_0^{9}\left(\sqrt{x}\right)^2dx=\pi\int_0^{9}1x\,dx=\pi\left[\dfrac{x^2}{2}\right]_0^{9}=\dfrac{81\pi}{2}\approx127{,}23$ m$^3$.}]{Một bồn chứa được tạo thành khi quay quanh trục hoành hình phẳng giới hạn bởi $y=\sqrt{x}$, $y=0$, $x=0$ và $x=9$ (đơn vị trên trục là mét). Thể tích của bồn chứa bằng bao nhiêu mét khối (làm tròn đến hàng phần trăm)?}
  {$127{,}23$.}
  {$254{,}47$.}{$40{,}5$.}{$63{,}62$.}

\cauTN[id=T12-C4-B13-D75-VD-011,muc=VD,loigiai={$V=\pi\int_0^{3}\left(\sqrt{4x}\right)^2dx=\pi\int_0^{3}4x\,dx=\pi\left[\dfrac{4x^2}{2}\right]_0^{3}=\dfrac{36\pi}{2}\approx56{,}55$ m$^3$.}]{Một bồn chứa được tạo thành khi quay quanh trục hoành hình phẳng giới hạn bởi $y=\sqrt{4x}$, $y=0$, $x=0$ và $x=3$ (đơn vị trên trục là mét). Thể tích của bồn chứa bằng bao nhiêu mét khối (làm tròn đến hàng phần trăm)?}
  {$56{,}55$.}
  {$113{,}1$.}{$18$.}{$28{,}27$.}

\cauTN[id=T12-C4-B13-D75-VD-012,muc=VD,loigiai={$V=\pi\int_0^{2}\left(\sqrt{4x}\right)^2dx=\pi\int_0^{2}4x\,dx=\pi\left[\dfrac{4x^2}{2}\right]_0^{2}=\dfrac{16\pi}{2}\approx25{,}13$ m$^3$.}]{Một bồn chứa được tạo thành khi quay quanh trục hoành hình phẳng giới hạn bởi $y=\sqrt{4x}$, $y=0$, $x=0$ và $x=2$ (đơn vị trên trục là mét). Thể tích của bồn chứa bằng bao nhiêu mét khối (làm tròn đến hàng phần trăm)?}
  {$25{,}13$.}
  {$50{,}27$.}{$8$.}{$12{,}57$.}

\cauDS[id=T12-C4-B13-D75-VD-013,muc=VD,loigiai={$V=\pi\int_0^{3}\left(\sqrt{5x}\right)^2dx=\pi\int_0^{3}5x\,dx=\pi\left[\dfrac{5x^2}{2}\right]_0^{3}=\dfrac{45\pi}{2}\approx70{,}69$ m$^3$.}]{Một bồn chứa được tạo thành khi quay quanh trục hoành hình phẳng giới hạn bởi $y=\sqrt{5x}$, $y=0$, $x=0$ và $x=3$ (đơn vị trên trục là mét).}{%
  \y{Thể tích được tính bởi $V=\pi\int_0^{3}5x\,dx$.}{d}
  \y{Thể tích bồn bằng $\dfrac{45\pi}{2}$ m$^3$.}{d}
  \y{Thể tích bồn xấp xỉ $70{,}69$ m$^3$.}{d}
  \y{Thể tích bồn bằng $45\pi$ m$^3$.}{s}}

\cauDS[id=T12-C4-B13-D75-VD-014,muc=VD,loigiai={$V=\pi\int_0^{9}\left(\sqrt{2x}\right)^2dx=\pi\int_0^{9}2x\,dx=\pi\left[\dfrac{2x^2}{2}\right]_0^{9}=\dfrac{162\pi}{2}\approx254{,}47$ m$^3$.}]{Một bồn chứa được tạo thành khi quay quanh trục hoành hình phẳng giới hạn bởi $y=\sqrt{2x}$, $y=0$, $x=0$ và $x=9$ (đơn vị trên trục là mét).}{%
  \y{Thể tích được tính bởi $V=\pi\int_0^{9}2x\,dx$.}{d}
  \y{Thể tích bồn bằng $\dfrac{162\pi}{2}$ m$^3$.}{d}
  \y{Thể tích bồn xấp xỉ $254{,}47$ m$^3$.}{d}
  \y{Thể tích bồn bằng $162\pi$ m$^3$.}{s}}

\cauDS[id=T12-C4-B13-D75-VD-015,muc=VD,loigiai={$V=\pi\int_0^{7}\left(\sqrt{x}\right)^2dx=\pi\int_0^{7}1x\,dx=\pi\left[\dfrac{x^2}{2}\right]_0^{7}=\dfrac{49\pi}{2}\approx76{,}97$ m$^3$.}]{Một bồn chứa được tạo thành khi quay quanh trục hoành hình phẳng giới hạn bởi $y=\sqrt{x}$, $y=0$, $x=0$ và $x=7$ (đơn vị trên trục là mét).}{%
  \y{Thể tích được tính bởi $V=\pi\int_0^{7}1x\,dx$.}{d}
  \y{Thể tích bồn bằng $\dfrac{49\pi}{2}$ m$^3$.}{d}
  \y{Thể tích bồn xấp xỉ $76{,}97$ m$^3$.}{d}
  \y{Thể tích bồn bằng $49\pi$ m$^3$.}{s}}

\cauDS[id=T12-C4-B13-D75-VD-016,muc=VD,loigiai={$V=\pi\int_0^{2}\left(\sqrt{5x}\right)^2dx=\pi\int_0^{2}5x\,dx=\pi\left[\dfrac{5x^2}{2}\right]_0^{2}=\dfrac{20\pi}{2}\approx31{,}42$ m$^3$.}]{Một bồn chứa được tạo thành khi quay quanh trục hoành hình phẳng giới hạn bởi $y=\sqrt{5x}$, $y=0$, $x=0$ và $x=2$ (đơn vị trên trục là mét).}{%
  \y{Thể tích được tính bởi $V=\pi\int_0^{2}5x\,dx$.}{d}
  \y{Thể tích bồn bằng $\dfrac{20\pi}{2}$ m$^3$.}{d}
  \y{Thể tích bồn xấp xỉ $31{,}42$ m$^3$.}{d}
  \y{Thể tích bồn bằng $20\pi$ m$^3$.}{s}}

\cauTLN[id=T12-C4-B13-D75-VD-017,muc=VD,loigiai={$V=\pi\int_0^{5}\left(\sqrt{5x}\right)^2dx=\pi\int_0^{5}5x\,dx=\pi\left[\dfrac{5x^2}{2}\right]_0^{5}=\dfrac{125\pi}{2}\approx196{,}35$ m$^3$.}]{Một bồn chứa được tạo thành khi quay quanh trục hoành hình phẳng giới hạn bởi $y=\sqrt{5x}$, $y=0$, $x=0$ và $x=5$ (đơn vị trên trục là mét). Tính thể tích bồn chứa (m$^3$, làm tròn đến hàng phần trăm).}{196,35}

\cauTLN[id=T12-C4-B13-D75-VD-018,muc=VD,loigiai={$V=\pi\int_0^{4}\left(\sqrt{x}\right)^2dx=\pi\int_0^{4}1x\,dx=\pi\left[\dfrac{x^2}{2}\right]_0^{4}=\dfrac{16\pi}{2}\approx25{,}13$ m$^3$.}]{Một bồn chứa được tạo thành khi quay quanh trục hoành hình phẳng giới hạn bởi $y=\sqrt{x}$, $y=0$, $x=0$ và $x=4$ (đơn vị trên trục là mét). Tính thể tích bồn chứa (m$^3$, làm tròn đến hàng phần trăm).}{25,13}

\cauTLN[id=T12-C4-B13-D75-VD-019,muc=VD,loigiai={$V=\pi\int_0^{4}\left(\sqrt{3x}\right)^2dx=\pi\int_0^{4}3x\,dx=\pi\left[\dfrac{3x^2}{2}\right]_0^{4}=\dfrac{48\pi}{2}\approx75{,}4$ m$^3$.}]{Một bồn chứa được tạo thành khi quay quanh trục hoành hình phẳng giới hạn bởi $y=\sqrt{3x}$, $y=0$, $x=0$ và $x=4$ (đơn vị trên trục là mét). Tính thể tích bồn chứa (m$^3$, làm tròn đến hàng phần trăm).}{75,4}

\cauTLN[id=T12-C4-B13-D75-VD-020,muc=VD,loigiai={$V=\pi\int_0^{4}\left(\sqrt{3x}\right)^2dx=\pi\int_0^{4}3x\,dx=\pi\left[\dfrac{3x^2}{2}\right]_0^{4}=\dfrac{48\pi}{2}\approx75{,}4$ m$^3$.}]{Một bồn chứa được tạo thành khi quay quanh trục hoành hình phẳng giới hạn bởi $y=\sqrt{3x}$, $y=0$, $x=0$ và $x=4$ (đơn vị trên trục là mét). Tính thể tích bồn chứa (m$^3$, làm tròn đến hàng phần trăm).}{75,4}

\cauTN[id=T12-C4-B13-D75-VDC-001,muc=VDC,loigiai={$V=\pi\int_0^{9}\left(\sqrt{4x}\right)^2dx=\pi\int_0^{9}4x\,dx=\pi\left[\dfrac{4x^2}{2}\right]_0^{9}=\dfrac{324\pi}{2}\approx508{,}94$ m$^3$.}]{Một bồn chứa được tạo thành khi quay quanh trục hoành hình phẳng giới hạn bởi $y=\sqrt{4x}$, $y=0$, $x=0$ và $x=9$ (đơn vị trên trục là mét). Thể tích của bồn chứa bằng bao nhiêu mét khối (làm tròn đến hàng phần trăm)?}
  {$508{,}94$.}
  {$1017{,}88$.}{$162$.}{$254{,}47$.}

\cauTN[id=T12-C4-B13-D75-VDC-002,muc=VDC,loigiai={$V=\pi\int_0^{6}\left(\sqrt{3x}\right)^2dx=\pi\int_0^{6}3x\,dx=\pi\left[\dfrac{3x^2}{2}\right]_0^{6}=\dfrac{108\pi}{2}\approx169{,}65$ m$^3$.}]{Một bồn chứa được tạo thành khi quay quanh trục hoành hình phẳng giới hạn bởi $y=\sqrt{3x}$, $y=0$, $x=0$ và $x=6$ (đơn vị trên trục là mét). Thể tích của bồn chứa bằng bao nhiêu mét khối (làm tròn đến hàng phần trăm)?}
  {$169{,}65$.}
  {$339{,}29$.}{$54$.}{$84{,}82$.}

\cauTN[id=T12-C4-B13-D75-VDC-003,muc=VDC,loigiai={$V=\pi\int_0^{3}\left(\sqrt{4x}\right)^2dx=\pi\int_0^{3}4x\,dx=\pi\left[\dfrac{4x^2}{2}\right]_0^{3}=\dfrac{36\pi}{2}\approx56{,}55$ m$^3$.}]{Một bồn chứa được tạo thành khi quay quanh trục hoành hình phẳng giới hạn bởi $y=\sqrt{4x}$, $y=0$, $x=0$ và $x=3$ (đơn vị trên trục là mét). Thể tích của bồn chứa bằng bao nhiêu mét khối (làm tròn đến hàng phần trăm)?}
  {$56{,}55$.}
  {$113{,}1$.}{$18$.}{$28{,}27$.}

\cauTN[id=T12-C4-B13-D75-VDC-004,muc=VDC,loigiai={$V=\pi\int_0^{8}\left(\sqrt{x}\right)^2dx=\pi\int_0^{8}1x\,dx=\pi\left[\dfrac{x^2}{2}\right]_0^{8}=\dfrac{64\pi}{2}\approx100{,}53$ m$^3$.}]{Một bồn chứa được tạo thành khi quay quanh trục hoành hình phẳng giới hạn bởi $y=\sqrt{x}$, $y=0$, $x=0$ và $x=8$ (đơn vị trên trục là mét). Thể tích của bồn chứa bằng bao nhiêu mét khối (làm tròn đến hàng phần trăm)?}
  {$100{,}53$.}
  {$201{,}06$.}{$32$.}{$50{,}27$.}

\cauTN[id=T12-C4-B13-D75-VDC-005,muc=VDC,loigiai={$V=\pi\int_0^{4}\left(\sqrt{3x}\right)^2dx=\pi\int_0^{4}3x\,dx=\pi\left[\dfrac{3x^2}{2}\right]_0^{4}=\dfrac{48\pi}{2}\approx75{,}4$ m$^3$.}]{Một bồn chứa được tạo thành khi quay quanh trục hoành hình phẳng giới hạn bởi $y=\sqrt{3x}$, $y=0$, $x=0$ và $x=4$ (đơn vị trên trục là mét). Thể tích của bồn chứa bằng bao nhiêu mét khối (làm tròn đến hàng phần trăm)?}
  {$75{,}4$.}
  {$150{,}8$.}{$24$.}{$37{,}7$.}

\cauTN[id=T12-C4-B13-D75-VDC-006,muc=VDC,loigiai={$V=\pi\int_0^{7}\left(\sqrt{3x}\right)^2dx=\pi\int_0^{7}3x\,dx=\pi\left[\dfrac{3x^2}{2}\right]_0^{7}=\dfrac{147\pi}{2}\approx230{,}91$ m$^3$.}]{Một bồn chứa được tạo thành khi quay quanh trục hoành hình phẳng giới hạn bởi $y=\sqrt{3x}$, $y=0$, $x=0$ và $x=7$ (đơn vị trên trục là mét). Thể tích của bồn chứa bằng bao nhiêu mét khối (làm tròn đến hàng phần trăm)?}
  {$230{,}91$.}
  {$461{,}81$.}{$73{,}5$.}{$115{,}45$.}

\cauDS[id=T12-C4-B13-D75-VDC-007,muc=VDC,loigiai={$V=\pi\int_0^{9}\left(\sqrt{x}\right)^2dx=\pi\int_0^{9}1x\,dx=\pi\left[\dfrac{x^2}{2}\right]_0^{9}=\dfrac{81\pi}{2}\approx127{,}23$ m$^3$.}]{Một bồn chứa được tạo thành khi quay quanh trục hoành hình phẳng giới hạn bởi $y=\sqrt{x}$, $y=0$, $x=0$ và $x=9$ (đơn vị trên trục là mét).}{%
  \y{Thể tích được tính bởi $V=\pi\int_0^{9}1x\,dx$.}{d}
  \y{Thể tích bồn bằng $\dfrac{81\pi}{2}$ m$^3$.}{d}
  \y{Thể tích bồn xấp xỉ $127{,}23$ m$^3$.}{d}
  \y{Thể tích bồn bằng $81\pi$ m$^3$.}{s}}

\cauDS[id=T12-C4-B13-D75-VDC-008,muc=VDC,loigiai={$V=\pi\int_0^{5}\left(\sqrt{3x}\right)^2dx=\pi\int_0^{5}3x\,dx=\pi\left[\dfrac{3x^2}{2}\right]_0^{5}=\dfrac{75\pi}{2}\approx117{,}81$ m$^3$.}]{Một bồn chứa được tạo thành khi quay quanh trục hoành hình phẳng giới hạn bởi $y=\sqrt{3x}$, $y=0$, $x=0$ và $x=5$ (đơn vị trên trục là mét).}{%
  \y{Thể tích được tính bởi $V=\pi\int_0^{5}3x\,dx$.}{d}
  \y{Thể tích bồn bằng $\dfrac{75\pi}{2}$ m$^3$.}{d}
  \y{Thể tích bồn xấp xỉ $117{,}81$ m$^3$.}{d}
  \y{Thể tích bồn bằng $75\pi$ m$^3$.}{s}}

\cauTLN[id=T12-C4-B13-D75-VDC-009,muc=VDC,loigiai={$V=\pi\int_0^{4}\left(\sqrt{2x}\right)^2dx=\pi\int_0^{4}2x\,dx=\pi\left[\dfrac{2x^2}{2}\right]_0^{4}=\dfrac{32\pi}{2}\approx50{,}27$ m$^3$.}]{Một bồn chứa được tạo thành khi quay quanh trục hoành hình phẳng giới hạn bởi $y=\sqrt{2x}$, $y=0$, $x=0$ và $x=4$ (đơn vị trên trục là mét). Tính thể tích bồn chứa (m$^3$, làm tròn đến hàng phần trăm).}{50,27}

\cauTLN[id=T12-C4-B13-D75-VDC-010,muc=VDC,loigiai={$V=\pi\int_0^{9}\left(\sqrt{x}\right)^2dx=\pi\int_0^{9}1x\,dx=\pi\left[\dfrac{x^2}{2}\right]_0^{9}=\dfrac{81\pi}{2}\approx127{,}23$ m$^3$.}]{Một bồn chứa được tạo thành khi quay quanh trục hoành hình phẳng giới hạn bởi $y=\sqrt{x}$, $y=0$, $x=0$ và $x=9$ (đơn vị trên trục là mét). Tính thể tích bồn chứa (m$^3$, làm tròn đến hàng phần trăm).}{127,23}
